\documentclass[a4paper, 11pt]{article}
\usepackage[slovak]{babel}
\usepackage[utf8]{inputenc}
\usepackage[left=2cm, top=3cm, text={17cm, 24cm}]{geometry}
\usepackage{times}
\usepackage[unicode]{hyperref}


\begin{document}

\begin{titlepage}
\begin{center}
\Huge
\textsc{Vysoké učení technické v Brně} \\ [0.3em]
\huge
\textsc{Fakulta informačních technologií} \\
\vspace{\stretch{0.382}}
\huge
IMP Projekt\\[0.2em]
\LARGE
Projektová Dokumentácia\\[0.1em]
\Large
\textbf{Meranie vzdialenosti ultrazvukovým senzorom}\\
\vspace{\stretch{0.618}}
\end{center}
{\Large 2025 \hfill
Jozef Sládeček (xsladej00)}
\end{titlepage}

\clearpage
\pagenumbering{arabic}

\newpage

\section{Úvod}

Projekt realizuje ultrazvukový merač vzdialenosti na platforme ESP32 s využitím ultrazvukového senzora HY-SRF05 a 4-ciferného 7-segmentového LED displeja. Vytvorený program umožňuje kontinuálne meranie vzdialenosti v rozsahu 0--400~cm a zobrazovanie nameraných hodnôt v reálnom čase. Používateľ môže pomocou tlačidiel ovládať základné funkcie zariadenia.

\section{Použité hardvérové a softvérové prostriedky}
\subsection{Hardvér}
\begin{itemize}
    \item \textbf{Vývojová doska Wemos D1 R32}
    \begin{itemize}
        \item Mikrokontrolér ESP32
    \end{itemize}
    
    \item \textbf{Ultrazvukový senzor HY-SRF05}
    \begin{itemize}
        \item Merací rozsah: 2--400~cm
        \item Presnosť: $\pm$3~mm
        \item Pripojenie: TRIG (GPIO 26), ECHO (GPIO 25)
    \end{itemize}
    
    \item \textbf{4-ciferný 7-segmentový LED displej s obvodom TLC5916}
    \begin{itemize}
        \item Ovládanie: SPI rozhranie (MOSI GPIO 23, SCLK GPIO 18)
        \item Riadenie multiplexovania: GPIO 14, 27, 16, 12
        \item PWM regulácia jasu: GPIO 13
    \end{itemize}
    
    \item \textbf{Ovládacie tlačidlá SW1--SW4}
    \begin{itemize}
        \item SW1 (GPIO 4): Prepínanie režimu zobrazenia
        \item SW2 (GPIO 5): Nastavenie jasu displeja
        \item SW3 (GPIO 33): Zmrazenie/odmrazenie hodnoty
        \item SW4 (GPIO 32): Správa histórie meraní
    \end{itemize}
\end{itemize}
\subsection{Softvér}
\begin{itemize}
    \item \textbf{ESP-IDF v6.1}: Vývojový framework
    \item \textbf{FreeRTOS}: RTOS (Real-Time Operating System)
    \begin{itemize}
        \item 3 paralelné tasky (Display, Measurement, Button)
        \item Všetky s prioritou 1
    \end{itemize}
\end{itemize}

\subsubsection{Kompilácia a nahratie programu}
Aktivácia ESP-IDF prostredia:
\begin{verbatim}
. ~/esp/esp-idf/export.sh
\end{verbatim}

Kompilácia, nahratie a spustenie monitora:
\begin{verbatim}
idf.py -p /dev/ttyUSB0 flash monitor
\end{verbatim}

\section{Návrh a realizácia riešenia}

\subsection{Architektúra systému}
Program je navrhnutý ako multitaskový systém s využitím FreeRTOS. Systém je rozdelený na tri nezávislé tasky:

\begin{itemize}
    \item \textbf{Display Task (priorita 1)}: Kontinuálne zobrazovanie aktuálnej hodnoty vzdialenosti na 4-cifernom displeji pomocou multiplexovania. Frekvencia obnovovania je približne 125~Hz (8~ms na kompletný cyklus 4 číslic).
    
    \item \textbf{Measurement Task (priorita 1)}: Periodické meranie vzdialenosti ultrazvukovým senzorom s intervalem 200~ms (5~Hz). Nameraná hodnota sa ukladá do zdieľanej premennej \texttt{current\_distance}.
    
    \item \textbf{Button Task (priorita 1)}: Obsluha tlačidiel s debouncing mechanizmom (200~ms). Spracováva prerušenia od tlačidiel SW1--SW3 a polling tlačidla SW4 pre detekciu krátkych/dlhých stlačení.
\end{itemize}

Tasky zdieľajú dáta prostredníctvom globálnych premenných označených ako \texttt{volatile}. Funkcia \texttt{app\_main()} inicializuje všetky komponenty a vytvára tasky pomocou \texttt{xTaskCreate()}. Po vytvorení taskov FreeRTOS scheduler prevezme kontrolu.

\subsection{Princíp merania vzdialenosti}
Ultrazvukový senzor HY-SRF05 využíva metódu meraného odrazu ultrazvukových vĺn:

\begin{enumerate}
    \item Mikrokontrolér vyšle na pin TRIG 10~$\mu$s impulz
    \item Senzor vyšle 8 ultrazvukových impulzov (40~kHz)
    \item Pin ECHO prejde do stavu HIGH
    \item Po prijatí odrazeného signálu prejde ECHO do stavu LOW
    \item Vzdialenost sa vypočíta podľa vzorca:
    \[
    d = \frac{t \cdot v}{2}
    \]
    kde $t$ je čas pulzu v $\mu$s, $v = 0.0343$~cm/$\mu$s (rýchlosť zvuku), delenie 2 kompenzuje cestu tam aj späť
\end{enumerate}

Funkcia \texttt{measure\_distance()} implementuje merací cyklus s timeout mechanizmom. Čas šírky ECHO pulzu sa meria pomocou \texttt{esp\_timer\_get\_time()}, ktorá poskytuje mikrosekundovú presnosť. Pre zabránenie nekonečného čakania je implementovaný timeout 30~ms (zodpovedá max. vzdialenosti ~5~m). Pri chybnom meraní alebo prekročení rozsahu funkcia vráti hodnotu -1, čo sa zobrazí ako \texttt{----}.

\subsection{Ovládanie displeja}
4-ciferný 7-segmentový displej je ovládaný pomocou multiplexovania a sériového prenosu dát:

\begin{itemize}
    \item \textbf{Prenos dát}: Segment pattern (8-bitová hodnota definujúca stav jednotlivých segmentov) sa prenáša bit po bite pomocou MOSI a SCLK pinov. Funkcia \texttt{send\_byte()} iteruje cez 8 bitov a generuje clock impulzy. Po prenesení celého bajtu sa dáta latchujú signálom XLAT pomocou funkcie \texttt{latch\_data()}
    
    \item \textbf{Multiplexovanie}: Každá číslica je aktívna 2~ms, potom sa prepne na nasledujúcu. Celý cyklus trvá 8~ms (125~Hz), čo eliminuje blikanie displeja. Funkcia \texttt{show\_digit()} vypne všetky cifry, prenesie segment pattern a aktivuje príslušnú číslicu
    
    \item \textbf{Segment patterns}: Pre každú cifru 0--9 je definovaný 8-bitový pattern (\texttt{DIGIT\_0} až \texttt{DIGIT\_9}), kde každý bit reprezentuje jeden segment
    
    \item \textbf{Implementácia režimov zobrazenia}: Funkcia \texttt{display\_distance()} implementuje tri režimy (STATIC/ADAPTIVE/INTEGER). Hodnota sa vynásobí faktorom 10 alebo 100, rozdelí na cifry pomocou \texttt{extract\_digits()} a desatinná bodka sa pridá na príslušnú pozíciu. Detailný popis jednotlivých režimov je v sekcii 4.2
    
    \item \textbf{PWM regulácia jasu}: Pin BLANK je riadený PWM signálom pomocou LEDC drivera. Jas je regulovaný pomocou duty cycle (invertovaná logika):
    \begin{itemize}
        \item HIGH: 0\% duty → 100\% jas
        \item MEDIUM: 50\% duty → 50\% jas
        \item LOW: 78\% duty → 22\% jas
        \item MIN: 90\% duty → 10\% jas
    \end{itemize}
\end{itemize}
\subsection{Obsluha tlačidiel}
Tlačidlá sú obsluhované kombináciou prerušení (ISR) a pollingu:

\begin{itemize}
    \item \textbf{SW1, SW2, SW3}: Využívajú GPIO prerušenia (NEGEDGE). ISR handler iba nastaví príslušný flag, samotné spracovanie prebieha v Button Task, ktorý kontroluje flagy každých 20~ms a implementuje debouncing 200~ms
    
    \item \textbf{SW4}: Používa polling s detekciou trvania stlačenia. Pri stlačení sa zaznamená čas pomocou \texttt{esp\_timer\_get\_time()}, pri uvolnení sa vyhodnotí trvanie (hranica 1~s pre rozlíšenie krátkého a dlhého stlačenia)
\end{itemize}

História meraní je implementovaná ako kruhový buffer s kapacitou 4 hodnoty (pole \texttt{measurement\_history[4]}). Pri zmrazení aktuálnej hodnoty (SW3) sa táto automaticky uloží pomocou funkcie \texttt{save\_to\_history()}, ktorá posunie staršie hodnoty a vloží novú na začiatok bufferu.


\section{Funkčnosť}

Táto sekcia popisuje ovládanie zariadenia z pohľadu používateľa.

\subsection{Ovládacie tlačidlá}

\begin{itemize}
    \item \textbf{SW1}: Prepínanie režimu zobrazenia
    \begin{itemize}
        \item Statický režim -- desatinná bodka vždy na rovnakom mieste
        \item Adaptívny režim -- automatické prispôsobenie formátu pre maximálne využitie displeja
        \item Celočíselný režim -- zobrazenie bez desatinnej bodky
    \end{itemize}
    
    \item \textbf{SW2}: Nastavenie jasu displeja (4 úrovne cyklicky: HIGH → MEDIUM → LOW → MIN → HIGH)
    
    \item \textbf{SW3}: Zmrazenie/odmrazenie aktuálnej hodnoty
    \begin{itemize}
        \item Pri zmrazení sa hodnota uloží do histórie
        \item Meranie pokračuje na pozadí
        \item Opätovným stlačením sa obnoví živé zobrazenie
    \end{itemize}
    
    \item \textbf{SW4}: Práca s históriou meraní
    \begin{itemize}
        \item Dlhé stlačenie ($>1$ s): Vstup/výstup z režimu prehliadania histórie
        \item Krátke stlačenie ($<1$ s): Listovanie medzi uloženými meraniami (iba v režime histórie)
    \end{itemize}
\end{itemize}

\subsection{Režimy zobrazenia vzdialenosti}

\subsubsection{Statický režim}
Desatinná bodka je vždy na druhej pozícii zľava. Príklady zobrazení:
\begin{itemize}
    \item Pre 5.3 cm → displej ukáže \texttt{~~5.3}
    \item Pre 25.7 cm → displej ukáže \texttt{~25.7}
    \item Pre 123.4 cm → displej ukáže \texttt{123.4}
\end{itemize}

\subsubsection{Adaptívny režim}
Formát sa automaticky prispôsobuje rozsahu hodnoty pre zobrazenie maximálneho počtu platných cifier:
\begin{itemize}
    \item Pre 5.34 cm → displej ukáže \texttt{~5.34} (dve desatinné miesta)
    \item Pre 25.67 cm → displej ukáže \texttt{25.67} (dve desatinné miesta)
    \item Pre 123.4 cm → displej ukáže \texttt{123.4} (jedno desatinné miesto)
    \item Pre 1234 cm → displej ukáže \texttt{1234} (bez desatinnej bodky)
\end{itemize}

\subsubsection{Celočíselný režim}
Zobrazujú sa len celé centimetre bez desatinnej bodky:
\begin{itemize}
    \item Pre 5.7 cm → displej ukáže \texttt{~~~5}
    \item Pre 25.3 cm → displej ukáže \texttt{~~25}
    \item Pre 123.8 cm → displej ukáže \texttt{~123}
\end{itemize}

\subsection{História meraní}

Zariadenie automaticky ukladá posledné 4 zmrazené merania. História funguje ako kruhový buffer -- pri uložení piatej hodnoty sa prepíše najstaršia.

\textbf{Postup práce s históriou:}
\begin{enumerate}
    \item Zmrazte hodnotu tlačidlom SW3 (automaticky sa uloží)
    \item Dlhým stlačením SW4 ($>1$ s) vstúpte do režimu histórie
    \item Krátkym stlačením SW4 listujte medzi uloženými hodnotami
    \item Dlhým stlačením SW4 ($>1$ s) sa vráťte do normálneho režimu
\end{enumerate}

V režime histórie nie je možné zmrazovať nové hodnoty (SW3 je neaktívne).

\subsection{Indikácia chybových stavov}

Pri chybnom meraní sa na displeji zobrazí \texttt{----}. Tento stav indikuje:
\begin{itemize}
    \item Objekt je mimo rozsahu senzora ($< 2$ cm alebo $> 400$ cm)
    \item Nedostatočný odraz od povrchu (príliš hladký alebo mäkký materiál)
    \item Timeout -- senzor nezachytil odrazenú vlnu do 30 ms
\end{itemize}


\section{Záver}

Projekt bol úspešne implementovaný do finálneho stavu. Program realizuje ultrazvukový merač vzdialenosti s kontinuálnym meraním, zobrazovaním na 4-cifernom displeji a ovládaním pomocou tlačidiel. Všetky požiadavky zadania boli splnené.

Funkčnosť zariadenia bola overená praktickým testovaním -- porovnaním nameraných hodnôt s referenčnými vzdialenosťami meranými pravítkom. Zariadenie meria konzistentne v pracovnom rozsahu senzora. Prevádzka systému a všetky implementované funkcie sú dokumentované v priloženom demonštračnom videu\footnote{\url{https://drive.google.com/file/d/1zJ1TuUQlUMzdUARkRMVFFidXT5KiZ0wT/view?usp=sharing}}.

Implementovaný multitaskový systém s využitím FreeRTOS zabezpečuje oddelenie jednotlivých častí programu -- displej, meranie a obsluha tlačidiel bežia paralelne v samostatných taskoch.


\end{document}
